\section{The entropy gain across a cut in Hastings' area law}
\label{sec:hastings_relent}

% Hastings, arXiv:0705.2024v4, law16.tex, lines 234--242 (eq. relent),
% 413--503 (its proof), and 916--925 (note added, route 1).

The doubling step of \cite{Hastings2007AreaLaw} rests on a lower bound for the
mutual information $I(A:B)=S(A)+S(B)-S(A\cup B)$ of the two intervals of $l$
sites beside a cut $c$. The approximate factorization of the ground-state
projector $P_0=|\Omega\rangle\langle\Omega|$ in
\cite[Lemma~1]{Hastings2007AreaLaw} yields the effect $E=O_B^\dagger O_B$
(note added, route~1) with large expectation $\operatorname{tr}(P_0 E)$ and
small product-marginal expectation $\operatorname{tr}(\sigma E)$, so relative
entropy monotonicity bounds $I(A:B)$ from below, with constants $2C_1$ and
$C_2=7$ as recorded in~\cite{gap:hastings07_relent_constants}.

\begin{theorem}[The entropy gain from the factorization, stage form]
    \label{thm:hastings_relent_projector}
    \lean{TNLean.MPS.AreaLaw.Hastings.relent_of_projectorFactorization}
    \leanok
    \uses{def:chain_model}
    Let $\Omega$ be a unit vector on the chain of $L$ sites with local
    dimension $q$, let $c,l\in\N$, $\xi'>0$, and $C_1\ge1$. Set
    $A=[c-l,c)\cap[0,L)$, $B=[c,c+l)\cap[0,L)$,
    $P_0=|\Omega\rangle\langle\Omega|$, $\epsilon=C_1e^{-l/\xi'}$, and
    $\sigma=\rho_{[0,c)}\otimes\rho_{[c,L)}$. Suppose commuting orthogonal
    projections $O_L$ on $[0,c)$ and $O_R$ on $[c,L)$ and an operator $O_B$ on
    $A\cup B$ with $\|O_B\|\le1$ satisfy
    \begin{align}
        \|O_BO_LO_R-P_0\|\le\epsilon,\qquad
        \operatorname{tr}(\sigma P_0)\le2\epsilon.
        \label{eq:hastings_relent_hyp}
    \end{align}
    Then, with $\epsilon'=2\epsilon=2C_1e^{-l/\xi'}$,
    \begin{align}
        S(A)+S(B)-S(A\cup B)\ge\bigl(1-2\epsilon'\bigr)\frac{l}{\xi'}
        -\ln(2C_1)-7.
        \label{eq:hastings_relent}
    \end{align}
    This is a stage of eq.~(relent) of~\cite{Hastings2007AreaLaw},
    lines 234--242, with the constants
    of~\cite{gap:hastings07_relent_constants}: the factorization of Lemma~1
    and the overlap bound in~(\ref{eq:hastings_relent_hyp}) are hypotheses.
\end{theorem}
\begin{proof}\leanok
    \uses{lem:chain_entropy_inequalities}
    By subadditivity (Lemma~\ref{lem:chain_entropy_inequalities}),
    $S(A)+S(B)-S(A\cup B)\ge0$, which gives~(\ref{eq:hastings_relent}) when
    $\epsilon'>1/68$ since the right-hand side is negative. Assume
    $\epsilon'\le1/68$ and set $E=O_B^\dagger O_B$ and $Z=O_LO_R$. From
    (\ref{eq:hastings_relent_hyp}),
    \begin{align}
        \operatorname{tr}(P_0 E)=\|O_B\Omega\|^2\ge(1-2\epsilon)^2\ge1-2\epsilon',
        \qquad
        \|ZEZ-P_0\|\le2\epsilon+\epsilon^2,\qquad
        \operatorname{tr}(\sigma Z)\ge1-2\epsilon^2,
        \notag
    \end{align}
    and Cauchy--Schwarz for $\sigma$ yields
    $x=\operatorname{tr}(\sigma E)\le36\epsilon'\le36/68$. Since $E$ is
    supported on $A\cup B$, $\operatorname{tr}(P_0 E)=\operatorname{tr}(\rho_{A\cup B}E)$
    and $x=\operatorname{tr}((\rho_A\otimes\rho_B)E)$, with $x>0$ because
    $\operatorname{supp}\rho_{A\cup B}\subseteq\operatorname{supp}\rho_A\otimes\operatorname{supp}\rho_B$
    and $\operatorname{tr}(\rho_{A\cup B}E)>0$. Monotonicity of relative
    entropy under $\{E,1-E\}$ and of the binary relative entropy $D_2(p\Vert q)$
    then gives
    \begin{align}
        S(A)+S(B)-S(A\cup B)
        \ge D_2\bigl(\operatorname{tr}(P_0 E)\,\big\Vert\,x\bigr)
        \ge D_2(1-2\epsilon'\Vert 36\epsilon')
        \ge\bigl((1-2\epsilon')\ln(1/\epsilon')-5\bigr)-1,
        \notag
    \end{align}
    which simplifies to~(\ref{eq:hastings_relent}).
\end{proof}

\begin{theorem}[The entropy gain across a cut of large entropy, stage form]
    \label{thm:hastings_relent_window}
    \lean{TNLean.MPS.AreaLaw.Hastings.relent_of_projectorFactorization_of_le_entropy}
    \leanok
    \uses{def:chain_model, def:hastings_bootstrap_remainder}
    Let $\Omega$, $c$, $l$, $\xi'$, $C_1$, $A$, and $B$ be as in
    Theorem~\ref{thm:hastings_relent_projector}, with $q\ge2$ and
    $c+l\le L$. Suppose that for every $m\in\N$ with $c+m\le L$ there exist
    $O_B,O_L,O_R$ as in Theorem~\ref{thm:hastings_relent_projector} with $O_B$
    supported on $[c-m,c+m)\cap[0,L)$ and
    $\|O_BO_LO_R-P_0\|\le C_1e^{-m/\xi'}$, and that the cut entropy
    $S_c=S([0,c))$ satisfies
    \begin{align}
        S_c\ge\bigl(l+\xi'\ln C_1\bigr)\ln q+F(\xi',q)
        \label{eq:hastings_relent_window_cond}
    \end{align}
    (equivalently, $l\le(S_c-F(\xi',q))/\ln q-\xi'\ln C_1$). This is the
    window assumption of~\cite[line 483]{Hastings2007AreaLaw} with $l$, $S_c$,
    and $q$ in place of $l_0\ge l$, $S_{\mathrm{cut}}\le S_c$, and $D$.
    Then~(\ref{eq:hastings_relent}) holds. This is a stage of eq.~(relent)
    of~\cite[lines 427--503]{Hastings2007AreaLaw} with the constants
    of~\cite{gap:hastings07_relent_constants}, taking the factorization of
    Lemma~1 at every scale $m$ as a hypothesis.
\end{theorem}
\begin{proof}\leanok
    \uses{thm:hastings_relent_projector, thm:hastings_bootstrap_chain_stage}
    Since $\sigma=\rho_{[0,c)}\otimes\rho_{[c,L)}$ is a mixture of pure product
    states across $c$ (each of Schmidt rank $k=1$), the overlap is
    $P=\operatorname{tr}(\sigma P_0)=\sum_iw_i|\langle\psi_i,\Omega\rangle|^2$.
    If $P>2C_1e^{-l/\xi'}$, then $2C_1^2/P<C_1e^{l/\xi'}$, so
    Theorem~\ref{thm:hastings_bootstrap_chain_stage} with $k=1$ gives
    \begin{align}
        S_c\le\xi'\ln(2C_1^2/P)\ln q+F(\xi',q)
        <\bigl(l+\xi'\ln C_1\bigr)\ln q+F(\xi',q),
        \notag
    \end{align}
    contradicting~(\ref{eq:hastings_relent_window_cond}). Thus
    $P\le2C_1e^{-l/\xi'}$, and Theorem~\ref{thm:hastings_relent_projector}
    yields~(\ref{eq:hastings_relent}).
\end{proof}

\begin{theorem}[The entropy gain across a cut of a gapped chain]
    \label{thm:hastings_relent}
    \lean{TNLean.MPS.AreaLaw.Hastings.exists_relent_of_isGappedGroundState}
    \leanok
    \uses{def:chain_model, def:hastings_bootstrap_remainder,
        thm:hastings_lemma1_endpoint}
    For every range $R\in\N$ there is a constant $K\ge0$ depending only on $R$
    with the following property. Let $\Omega$ be a gapped ground vector with
    gap $\Delta E>0$ of a range-$R$ Hamiltonian on the chain of $L$ sites with
    local dimension $q$ and term norms at most $J\ge0$. Let $\tilde\xi$ be the
    decay length of Theorem~\ref{thm:hastings_lemma1_endpoint}, and set
    $\xi'=6\tilde\xi$ and $C_1=K\tilde\xi^3$. For every cut $c$ and $l\in\N$
    with $c+l\le L$, set $A=[c-l,c)\cap[0,L)$ and $B=[c,c+l)$. If $S_c$
    satisfies~(\ref{eq:hastings_relent_window_cond}), then~(\ref{eq:hastings_relent})
    holds with $\epsilon'=2C_1e^{-l/\xi'}$.

    This is eq.~(relent) of~\cite[lines 234--242]{Hastings2007AreaLaw} under
    the window assumption~(\ref{eq:hastings_relent_window_cond}) of line 483,
    with the constants of~\cite{gap:hastings07_relent_constants}, the printed
    rate $\xi'=6\tilde\xi$ (lines 141--142), and polynomial prefactor
    $C_1=K\tilde\xi^3$ (lines 258--262). Here $\tilde\xi$ has the form
    $\max(2v/\Delta E,\xi_C)$ from a chain Lieb--Robinson estimate with
    $J$-dependent prefactor, single-site second operator, and raised velocity,
    rather than a Lieb--Robinson pair $(v,\xi_C)$ in the sense of lines
    96--112~\cite{gap:hastings07_lemma1_decay_length,
        gap:hastings07_gaussian_average_exponent}.
\end{theorem}
\begin{proof}\leanok
    \uses{thm:hastings_relent_window, thm:hastings_lemma1_endpoint,
        thm:hastings_factorization_mono, lem:chain_entropy_inequalities}
    Theorem~\ref{thm:hastings_lemma1_endpoint} supplies $K_0\ge0$ and operators
    $O_B,O_L,O_R$ at every $(c,m)$ with
    $\|O_BO_LO_R-P_0\|\le K_0\tilde\xi^3e^{-m/(6\tilde\xi)}$. Setting
    $K=K_0+1$ and using $\tilde\xi\ge8$ gives $C_1=K\tilde\xi^3\ge1$ and
    $\|O_BO_LO_R-P_0\|\le C_1e^{-m/\xi'}$ by
    Theorem~\ref{thm:hastings_factorization_mono}. For $q\ge2$,
    Theorem~\ref{thm:hastings_relent_window} yields~(\ref{eq:hastings_relent});
    for $q=1$, $S_c=0<1+\ln(\xi'/2+1)=F(\xi',1)$, so
    (\ref{eq:hastings_relent_window_cond}) is vacuous.
\end{proof}
